% CP-VI-0093
% target_chapter: VI/39 — Weil Divisors
% solution_provenance: FRESH_CANONICAL_AUTHORED_SOLUTION

\subsection*{CP-VI-0093: A principal Weil divisor}

Let $X$ be an integral normal Noetherian scheme and $f\in K(X)^\times$. Define the principal Weil divisor $\operatorname{div}(f)$ and explain why only finitely many codimension-one points occur on an affine open where $f=a/b$.

\medskip
\noindent\textbf{Solution.}

For each codimension-one point $x$, normality makes $\mathcal O_{X,x}$ a DVR, with valuation $v_x$. Define
\[
\operatorname{div}(f)=\sum_{x\in X^{(1)}}v_x(f)[\overline{\{x\}}].
\]
On an affine Noetherian open where $f=a/b$, a nonzero valuation can occur only along codimension-one components of $V(a)$ or $V(b)$. Noetherianity gives finitely many irreducible components, so the sum is locally finite and, for a rational function, globally finite on a suitable finite affine cover.
